%%%%%%%%%%%%%%%%%%%%%%%%%%%%%%%%%%%%%%%%%%%%%%%%%%%%%%%%%%%%%%%%%%%%%%%
%%  Einheitlicher Stil fuer alle selbst erzeugten Abbildungen
%%
%%  Grundgedanke: Alle TikZ-Grafiken der Arbeit verwenden ausschliesslich
%%  RWTH-Blau, RWTH-Grau (= Schwarz 75 %) und Schwarz. Statt in den
%%  einzelnen Abbildungen konkrete Farben zu nennen, wird hier eine
%%  BEDEUTUNGSEBENE definiert: jede Rolle bekommt einen Namen, die
%%  Abbildungen benennen nur noch die Rolle. Aendert man eine Zuordnung
%%  hier, zieht die Aenderung durch das gesamte Dokument.
%%
%%  Zwei Ebenen:
%%    neutral (Grau)  - Rahmen, Gehaeuse, Struktur, unbeteiligte Teile
%%    aktiv   (Blau)  - Fokus der Abbildung, Signal, bewegte oder
%%                      bestromte Elemente, das Werkstueck
%%
%%  Graustufen-Druck: Blau 100 (Grauwert ~78) und Grau (~101) liegen im
%%  Schwarzweissdruck dicht beieinander. Eine Unterscheidung wird daher
%%  NIE allein ueber den Farbton getragen, sondern immer zusaetzlich
%%  ueber Linienstaerke oder Flaechenhelligkeit. Die unten gewaehlten
%%  Flaechen sind dafuer bewusst weit auseinandergelegt:
%%    figFlaeche  ECEDED (~237) | figTeilFlaeche 8EBAE5 (~180)
%%    figAkzentFlaeche C7DDF2 (~221) wird nur gegen Weiss gesetzt.
%%
%%  Benoetigt: styles/rwthcolors.tex muss vorher geladen sein.
%%%%%%%%%%%%%%%%%%%%%%%%%%%%%%%%%%%%%%%%%%%%%%%%%%%%%%%%%%%%%%%%%%%%%%%

%% --- neutrale Ebene: Struktur ----------------------------------------
\colorlet{figKontur}{RWTHSchwarz75}        % 646567  Standardkontur
\colorlet{figKonturHell}{RWTHSchwarz50}    % 9C9E9F  Hilfs- und Hinterlinien
\colorlet{figFlaeche}{RWTHSchwarz10}       % ECEDED  neutrale Flaeche
\colorlet{figFlaecheTief}{RWTHSchwarz25}   % CFD1D2  neutrale Flaeche, kraeftiger
\colorlet{figText}{RWTHSchwarz75}          % 646567  Beschriftung

%% --- aktive Ebene: Fokus ---------------------------------------------
\colorlet{figAkzent}{RWTHBlau100}          % 00549F  Hauptakzent
\colorlet{figAkzentMittel}{RWTHBlau75}     % 407FB7  zweite Stufe
\colorlet{figAkzentFlaeche}{RWTHBlau25}    % C7DDF2  Akzentflaeche
\colorlet{figAkzentZart}{RWTHBlau10}       % E8F1FA  Akzentflaeche, sehr hell

%% --- Werkstueck: dokumentweit identisch ------------------------------
%% Der Hairpin bzw. das gefoerderte Bauteil sieht in JEDER Abbildung
%% gleich aus. Das ist der optische rote Faden durch die Arbeit.
\colorlet{figTeilKontur}{RWTHBlau100}
\colorlet{figTeilFlaeche}{RWTHBlau50}      % 8EBAE5, im S/W-Druck klar dunkler

%% --- Linienstaerken ---------------------------------------------------
%% Zentral gesetzt, damit alle Abbildungen dieselbe Strichhierarchie
%% haben. Die Akzentstaerke ist bewusst deutlich groesser, damit die
%% Unterscheidung auch ohne Farbe traegt.
\newlength{\figLWfein}   \setlength{\figLWfein}{0.4pt}
\newlength{\figLWnormal} \setlength{\figLWnormal}{0.6pt}
\newlength{\figLWakzent} \setlength{\figLWakzent}{1.0pt}

%% --- gemeinsame TikZ-Stile -------------------------------------------
%% Werden ueber \tikzset global bereitgestellt und koennen in jeder
%% Abbildung direkt verwendet oder lokal ueberschrieben werden.
\tikzset{
  %% Flaechen und Konturen
  figstruktur/.style   = {draw=figKontur, line width=\figLWnormal, fill=figFlaeche},
  figstrukturleer/.style = {draw=figKontur, line width=\figLWnormal, fill=white},
  figaktiv/.style      = {draw=figAkzent, line width=\figLWakzent, fill=figAkzentFlaeche},
  figteil/.style       = {draw=figTeilKontur, line width=\figLWfein, fill=figTeilFlaeche},
  %% Beschriftung
  figlbl/.style        = {font=\sffamily\scriptsize, text=figText},
  figlblakzent/.style  = {font=\sffamily\scriptsize, text=figAkzent},
  %% Pfeile
  figpfeil/.style      = {-{Stealth[length=2mm,width=1.6mm]}, draw=figKontur,
                          line width=\figLWnormal},
  figpfeilakzent/.style= {-{Stealth[length=2mm,width=1.6mm]}, draw=figAkzent,
                          line width=\figLWakzent},
  figpfeilfein/.style  = {-{Stealth[length=1.6mm,width=1.3mm]}, draw=figKonturHell,
                          line width=\figLWfein},
  figmass/.style       = {{Stealth[length=1.4mm,width=1.1mm]}-{Stealth[length=1.4mm,width=1.1mm]},
                          draw=figAkzent, line width=\figLWfein},
}
